\documentclass[11pt]{article}

\usepackage[a4paper,margin=27mm]{geometry}
\usepackage[T1]{fontenc}
\usepackage[utf8]{inputenc}
\usepackage[english]{babel}
\usepackage{amsmath,amssymb,amsfonts,amsthm,mathtools}
\usepackage{bm}
\usepackage{booktabs}
\usepackage{microtype}
\usepackage[numbers,sort&compress]{natbib}
\usepackage[colorlinks=true,citecolor=blue,linkcolor=magenta,urlcolor=blue]{hyperref}
\hypersetup{
  pdftitle={A no-go theorem for non-multiplexed conditional IQP universality},
  pdfsubject={Supplementary theoretical material}
}

\newtheorem{theorem}{Theorem}[section]
\newtheorem{proposition}[theorem]{Proposition}
\newtheorem{lemma}[theorem]{Lemma}
\newtheorem{corollary}[theorem]{Corollary}
\theoremstyle{definition}
\newtheorem{definition}[theorem]{Definition}
\theoremstyle{remark}
\newtheorem{remark}[theorem]{Remark}

\newcommand{\bits}{\{0,1\}}
\newcommand{\TVD}{\operatorname{TVD}}
\newcommand{\Tr}{\operatorname{Tr}}
\newcommand{\supp}{\operatorname{supp}}
\newcommand{\ketbra}[2]{|#1\rangle\!\langle#2|}
\newcommand{\braket}[2]{\langle#1|#2\rangle}
\newcommand{\cH}{\mathcal{H}}
\newcommand{\cM}{\mathcal{M}}
\newcommand{\cC}{\mathcal{C}}
\newcommand{\cP}{\mathcal{P}}
\newcommand{\R}{\mathbb{R}}

\title{A no-go theorem for non-multiplexed universality of\\
state-conditioned IQP circuits\\[2mm]
\large Supplementary theoretical material}
\author{}
\date{27 August 2026}

\begin{document}
\maketitle

\begin{abstract}
We ask whether a single instantaneous quantum polynomial-time (IQP) model,
whose circuit parameters are shared across conditions and whose condition
appears only in a classically tractable input state, can be a universal
conditional generator without using the condition as a selector for a bank of
unconditional generators.  It cannot.  The obstruction is independent of the
IQP gate set.  Any fixed circuit followed by measurement is a
quantum-to-classical channel, so the total-variation distance between two
output laws cannot exceed the trace distance between their input states.
Consequently, exact universality over conditional laws requires the input
states for every two distinct conditions to have orthogonal supports.  Uniform
approximation of all continuous conditional laws by a resource-scaled family
requires the corresponding states to become asymptotically perfectly
distinguishable.  Perfectly distinguishable condition states are operational
labels and permit measurement followed by branch selection; this is precisely
the routing resource excluded by a non-multiplexed notion of universality.

We also prove that a condition register which controls an IQP diagonal and is
then discarded produces a classical mixture of condition-independent IQP
branches.  Thus selector and Bernstein constructions do not evade the
obstruction; they implement multiplexing in coherent notation.  For the
paper's broadcast product encoder, we derive an explicit lower bound on the
uniform approximation error.  A target that changes from one deterministic
visible string to another over an interval of width $h$ has error at least
$\tfrac12(1-\sqrt{N}L_\alpha h/2)_+$ for $N$ encoding qubits and an
$L_\alpha$-Lipschitz angle.  At the reported Gaussian setting
$N=19$, $L_\alpha=0.70\pi$, and $h=0.01$, this lower bound is approximately
$0.476$.  Fixed-width Fourier and parameter-dimension obstructions give
separate counterexamples for the tested complete two-local models.  Classical
tractability of the inputs changes estimator cost, not these expressivity
bounds.
\end{abstract}

\section{Question and answer}
\label{sec:answer}

The model in the accompanying paper uses one circuit for all conditions.  A
condition $c$ changes an input state $\rho_c$; a fixed IQP unitary, hidden-qubit
discarding, and computational-basis measurement produce a visible law
$q_\theta(\cdot\mid c)$.  The question addressed here is deliberately
stronger than whether several unconditional IQP generators can be placed in
orthogonal blocks and selected by $c$.

\begin{definition}[Label separation and non-multiplexed universality]
Let $\cC$ be a condition space and let $\rho_c^{(R)}$ denote the condition
state at resource size $R$.  The encoder family is \emph{asymptotically
label-separating} if
\begin{equation}
  \sup_R D_{\rm tr}\!\left(\rho_c^{(R)},\rho_{c'}^{(R)}\right)=1
  \quad\text{for every }c\neq c'.
  \label{eq:label-separating-definition}
\end{equation}
We call it \emph{non-multiplexed} if it is not label-separating, so at least
one pair retains a trace-distance gap from one uniformly over resources.  A
model family is \emph{universal} for a target class $\cP$ if, for every
$p\in\cP$ and every $\varepsilon>0$, some resource size and some shared
circuit parameters satisfy
\begin{equation}
  \sup_{c\in\cC}\TVD\!\left(q_\theta(\cdot\mid c),
  p(\cdot\mid c)\right)<\varepsilon.
  \label{eq:uniform-universality}
\end{equation}
\end{definition}

The label-separation criterion formalizes the requested exclusion of ``stitching
together'' condition-addressed IQP circuits.  Sections~\ref{sec:contraction}
and~\ref{sec:resource} prove that universality for all conditional laws, or
even for all total-variation-continuous laws on $[0,1]$, contradicts this
clause.  The result is a counterexample to the broad universality claim:

\begin{quote}
\emph{Classically tractable conditional inputs do not make a shared IQP
architecture conditionally universal.  Exact universality forces orthogonal
condition labels; asymptotic universality for all continuous kernels forces
asymptotically orthogonal labels.  Excluding that routing limit excludes the
claimed universality.}
\end{quote}

The conclusion does not rule out approximation of a narrower target class
compatible with the encoding.  It identifies a necessary compatibility inequality
and quantifies its violation.

\section{Model and distances}
\label{sec:model}

Let $V$ contain $n\geq1$ visible qubits, let $H$ contain any number of hidden
qubits, and let $E$ contain any ancillary or condition-encoding qubits.  The
input state at condition $c\in\cC$ is a density operator $\rho_c$ on
$VHE$.  A fixed parameter choice $\theta$ specifies all subsequent operations:
an IQP unitary, discarding of $HE$, and measurement of $V$.  These operations
define a quantum-to-classical channel
\begin{equation}
  \cM_\theta(\rho)(x)=\Tr(M_x^{(\theta)}\rho),
  \qquad x\in\bits^n,
  \label{eq:povm-channel}
\end{equation}
where $\{M_x^{(\theta)}\}_x$ is the input-picture positive-operator-valued
measure (POVM).  Hence
\begin{equation}
  q_\theta(x\mid c)=\cM_\theta(\rho_c)(x).
  \label{eq:conditional-output}
\end{equation}
Nothing in the first no-go theorem depends on the IQP form of the channel.
It therefore holds for arbitrary diagonals, arbitrary hidden width, and
arbitrary trainable phase locality.

We take the encoding map to be fixed at each resource size, as in the paper.
If the encoder has shared trainable parameters $\varphi$, one may replace the
resource index $R$ everywhere below by the pair $(R,\varphi)$.  All proofs are
unchanged; the necessary suprema then range over every admissible encoder.

For probability laws $p,q$ on a finite set and density operators $\rho,\sigma$,
define
\begin{align}
  \TVD(p,q)&=\frac12\sum_x|p(x)-q(x)|,
  \label{eq:tvd-definition}\\
  D_{\rm tr}(\rho,\sigma)&=\frac12\lVert\rho-\sigma\rVert_1.
  \label{eq:trace-definition}
\end{align}
We use $(a)_+=\max\{a,0\}$.  These distance facts are standard; a
self-contained reference is \citet{watrousTheoryQuantumInformation2018}.
Total variation also has the event
characterization
\begin{equation}
  \TVD(p,q)=\max_A |p(A)-q(A)|.
  \label{eq:tvd-event}
\end{equation}

\section{The contraction obstruction}
\label{sec:contraction}

\begin{lemma}[Measurement contracts distinguishability]
\label{lem:contraction}
For every quantum-to-classical channel $\cM$ and density operators
$\rho,\sigma$,
\begin{equation}
  \TVD\bigl(\cM(\rho),\cM(\sigma)\bigr)
  \leq D_{\rm tr}(\rho,\sigma).
  \label{eq:contraction}
\end{equation}
\end{lemma}

\begin{proof}
Write $\cM(\rho)(x)=\Tr(M_x\rho)$.  For an output event $A$, let
$M_A=\sum_{x\in A}M_x$.  The POVM conditions give $0\leq M_A\leq I$.
The variational characterization of trace distance is
\begin{equation}
 D_{\rm tr}(\rho,\sigma)
 =\max_{0\leq F\leq I}|\Tr[F(\rho-\sigma)]|.
 \label{eq:trace-variational}
\end{equation}
Using Eq.~\eqref{eq:tvd-event},
\begin{align}
 \TVD\bigl(\cM(\rho),\cM(\sigma)\bigr)
 &=\max_A |\Tr[M_A(\rho-\sigma)]| \\
 &\leq\max_{0\leq F\leq I}|\Tr[F(\rho-\sigma)]|,
\end{align}
which is Eq.~\eqref{eq:contraction}.
\end{proof}

The following inequality is the central quantitative result.

\begin{theorem}[Pairwise approximation lower bound]
\label{thm:pairwise}
For any two conditions $c,c'$, any target conditional law $p$, and any fixed
channel $\cM_\theta$,
\begin{multline}
 \max\!\left\{
  \TVD\bigl(q_\theta(\cdot\mid c),p(\cdot\mid c)\bigr),
  \TVD\bigl(q_\theta(\cdot\mid c'),p(\cdot\mid c')\bigr)
 \right\}\\
 \geq \frac12\left[
  \TVD\bigl(p(\cdot\mid c),p(\cdot\mid c')\bigr)
  -D_{\rm tr}(\rho_c,\rho_{c'})
 \right]_+.
 \label{eq:pairwise-lower-bound}
\end{multline}
\end{theorem}

\begin{proof}
Set $e_c=\TVD(q_\theta(\cdot\mid c),p(\cdot\mid c))$, and similarly
define $e_{c'}$.  Two applications of the triangle inequality and
Lemma~\ref{lem:contraction} give
\begin{align}
 \TVD\bigl(p(\cdot\mid c),p(\cdot\mid c')\bigr)
 &\leq e_c+
 \TVD\bigl(q_\theta(\cdot\mid c),q_\theta(\cdot\mid c')\bigr)+e_{c'}\\
 &\leq e_c+D_{\rm tr}(\rho_c,\rho_{c'})+e_{c'}\\
 &\leq 2\max\{e_c,e_{c'}\}+D_{\rm tr}(\rho_c,\rho_{c'}).
\end{align}
Rearranging, and using nonnegativity of the left-hand side of
Eq.~\eqref{eq:pairwise-lower-bound}, proves the claim.
\end{proof}

Equation~\eqref{eq:pairwise-lower-bound} is an expressivity certificate for
any proposed conditional encoding.  A target can separate two conditions no
more sharply than their encoded states, up to twice the approximation error.

\begin{corollary}[Exact universality requires orthogonal labels]
\label{cor:orthogonal}
Suppose a fixed input encoding and a trainable family of channels represent
every conditional law on a condition set containing two distinct points.
Then
\begin{equation}
  D_{\rm tr}(\rho_c,\rho_{c'})=1
  \qquad\text{for every }c\neq c'.
  \label{eq:orthogonal-necessary}
\end{equation}
Equivalently, $\rho_c$ and $\rho_{c'}$ have orthogonal supports.
\end{corollary}

\begin{proof}
For a selected pair $c\neq c'$, choose a target with
$p(\cdot\mid c)=\delta_{0^n}$ and
$p(\cdot\mid c')=\delta_{1^n}$.  These laws have total-variation distance
one.  Exact representation makes the left-hand side of
Eq.~\eqref{eq:pairwise-lower-bound} zero, so
$D_{\rm tr}(\rho_c,\rho_{c'})\geq1$.  Trace distance between states is at most
one, hence equality holds.  Equality is equivalent to orthogonality of their
supports.  Repeating the argument for every pair gives the result.
\end{proof}

Pairwise orthogonal states can be distinguished with zero error by one
measurement.  A model supplied with such states therefore receives a
classical condition label in operational form: it can measure the label and
route computation to a condition-specific branch.  Exact universality does
not produce a genuinely shared conditional mechanism; it requires the
multiplexing resource that the question excludes.

\begin{corollary}[A classically tractable counterexample]
If $\rho_c=\rho_0$ for every $c$, then
$q_\theta(\cdot\mid c)=q_\theta(\cdot\mid c')$ for every pair of conditions
and every parameter choice.  The family cannot represent any target that
differs between two conditions.
\end{corollary}

The constant state may be a computational-basis product state.  It is
classically tractable under the strongest standard notions: it can be sampled
exactly and its amplitudes can be evaluated exactly.  Thus ``classically
tractable input states'' alone is not a sufficient premise for conditional
universality.

\section{Resource scaling does not remove the obstruction}
\label{sec:resource}

An asymptotic universality statement may choose a new input width and circuit
for each requested accuracy.  Index these resources by $R$ and write the
condition state as $\rho_c^{(R)}$.

\begin{theorem}[Necessary asymptotic separation]
\label{thm:asymptotic-separation}
Let $\cC=[0,1]$, and suppose a resource-scaled family is universal in the
sense of Eq.~\eqref{eq:uniform-universality} for every conditional law
$p(\cdot\mid c)$ that is continuous in total variation.  Then, for every
$c\neq c'$,
\begin{equation}
  \sup_R D_{\rm tr}\!\left(\rho_c^{(R)},\rho_{c'}^{(R)}\right)=1.
  \label{eq:asymptotic-separation}
\end{equation}
More precisely, approximation with uniform error $\varepsilon$ of a suitable
continuous target requires
\begin{equation}
  D_{\rm tr}\!\left(\rho_c^{(R)},\rho_{c'}^{(R)}\right)
  \geq 1-2\varepsilon.
  \label{eq:epsilon-separation}
\end{equation}
\end{theorem}

\begin{proof}
Assume without loss of generality that $c<c'$.  Define a continuous function
$\lambda:[0,1]\to[0,1]$ that equals zero on $[0,c]$, is linear from zero to
one on $[c,c']$, and equals one on $[c',1]$.  The target
\begin{equation}
 p(\cdot\mid t)=(1-\lambda(t))\delta_{0^n}
                  +\lambda(t)\delta_{1^n}
 \label{eq:ramp-target}
\end{equation}
is continuous in total variation and its endpoint laws have distance one.
Apply Theorem~\ref{thm:pairwise} to a resource $R$ that approximates this
target with uniform error $\varepsilon$.  Rearrangement yields
Eq.~\eqref{eq:epsilon-separation}.  Sending $\varepsilon$ to zero gives
Eq.~\eqref{eq:asymptotic-separation}.
\end{proof}

The theorem permits a different resource $R$ for different targets and
accuracies.  It still forces every fixed pair of distinct conditions to become
arbitrarily distinguishable somewhere in the family.

\begin{definition}[Uniform overlap preservation]
A resource-scaled encoder is uniformly overlap preserving if either of the
following sufficient conditions holds:
\begin{enumerate}
  \item there are $c\neq c'$ and $\gamma>0$ such that
  $D_{\rm tr}(\rho_c^{(R)},\rho_{c'}^{(R)})\leq1-\gamma$ for every $R$; or
  \item there is a common modulus $\omega$ with $\omega(\delta)\to0$ as
  $\delta\to0$ and
  \begin{equation}
    D_{\rm tr}(\rho_c^{(R)},\rho_{c'}^{(R)})
    \leq\omega(|c-c'|)
    \quad\text{for every }R,c,c'.
    \label{eq:uniform-modulus}
  \end{equation}
\end{enumerate}
\end{definition}

\begin{corollary}[No-go for overlap-preserving families]
\label{cor:overlap-no-go}
Under condition 1, the ramp target in Eq.~\eqref{eq:ramp-target} has uniform
approximation error at least $\gamma/2$.  Under condition 2, there is a
total-variation-continuous ramp target whose uniform approximation error is
at least $1/4$ for every resource size and every circuit.
\end{corollary}

\begin{proof}
The first statement follows immediately from
Eq.~\eqref{eq:pairwise-lower-bound}.  For the second, choose $h>0$ with
$\omega(h)\leq1/2$, take two points at distance $h$, and use the continuous
ramp between them.  Theorem~\ref{thm:pairwise} gives
$\tfrac12(1-\omega(h))\geq1/4$.
\end{proof}

Therefore, a universal resource-scaled family must lose every continuity
bound of the form Eq.~\eqref{eq:uniform-modulus} that is uniform over resource
size.  Its condition encoding must sharpen with resources until nearby conditions act as
nearly orthogonal labels.  This is label amplification, not a shared smooth
conditional representation.

\section{Why discarded IQP selectors are classical multiplexers}
\label{sec:selectors}

One might try to evade the preceding conclusion by preparing a selector in
superposition, applying a single diagonal unitary, and discarding the selector.
The next proposition shows that this notation implements a classical mixture
of branches.

Let $S$ be a selector register and $R$ collect all retained and hidden model
registers.  Consider the state immediately before the IQP diagonal.  Assume
that it factorizes as $\sigma_S(c)\otimes\tau_R$ and that only
$\sigma_S(c)$ depends on $c$.  Every unitary diagonal in the joint
computational basis has the controlled form
\begin{equation}
 D=\sum_{s\in\bits^r}\ketbra{s}{s}_S\otimes D_s,
 \label{eq:controlled-diagonal}
\end{equation}
where each $D_s$ is diagonal on $R$.

\begin{proposition}[Discarded-selector mixture identity]
\label{prop:selector-mixture}
After applying $D$ and discarding $S$, the state on $R$ is
\begin{equation}
 \Tr_S\!\left[D(\sigma_S(c)\otimes\tau_R)D^\dagger\right]
 =\sum_{s\in\bits^r} w_c(s)D_s\tau_R D_s^\dagger,
 \qquad
 w_c(s)=\langle s|\sigma_S(c)|s\rangle.
 \label{eq:selector-mixture-state}
\end{equation}
Consequently, after any fixed channel on $R$, with $S$ acted on only locally
and then discarded, the visible measurement obeys
\begin{equation}
 q(x\mid c)=\sum_s w_c(s)q_s(x),
 \label{eq:selector-mixture-law}
\end{equation}
where $q_s$ is a condition-independent branch law.  Off-diagonal coherence
of $\sigma_S(c)$ has no effect on the output.
\end{proposition}

\begin{proof}
Expand $\sigma_S(c)=\sum_{s,t}\sigma_{st}(c)\ketbra{s}{t}$.  Equation
\eqref{eq:controlled-diagonal} gives
\begin{equation}
 D(\sigma_S(c)\otimes\tau_R)D^\dagger
 =\sum_{s,t}\sigma_{st}(c)\ketbra{s}{t}
   \otimes D_s\tau_R D_t^\dagger.
\end{equation}
Since $\Tr_S(\ketbra{s}{t})=\delta_{st}$, only terms with $s=t$ survive,
which proves Eq.~\eqref{eq:selector-mixture-state}.  Linearity of the remaining
channel proves Eq.~\eqref{eq:selector-mixture-law}.  A final unitary acting
only on $S$, including the final IQP Hadamards, cannot change its partial
trace and therefore does not alter the identity.
\end{proof}

The identity covers finite one-hot selectors, binary-address registers, and
Bernoulli selectors whose weights form Bernstein polynomials.  Such
constructions can be written as one large diagonal, but they remain a bank of
$q_s$ distributions joined by condition-dependent classical weights.  Exact
branch selection further requires distinguishable selector states, in
agreement with Corollary~\ref{cor:orthogonal}.  These constructions are valid
compilations, but they are not a general non-multiplexed solution.

The factorization assumption is essential to this structural identity and is
satisfied by the usual product selector constructions; the first IQP
Hadamards preserve it.  An encoder that entangles $S$ and $R$ need not obey
Eq.~\eqref{eq:selector-mixture-law}, but it still obeys the contraction bound
in Theorem~\ref{thm:pairwise}.

\section{Specialization to product rotation encoders}
\label{sec:product}

The paper uses tractable product inputs generated by $Y$ rotations.  Consider
the general form
\begin{equation}
 |\Phi(c)\rangle
 =\bigotimes_{j=1}^{N}R_Y(\alpha_j(c))|0\rangle,
 \qquad
 R_Y(a)|0\rangle=\cos(a/2)|0\rangle+\sin(a/2)|1\rangle.
 \label{eq:product-encoder}
\end{equation}
The shared first Hadamard layer of an IQP circuit does not change pairwise
overlaps.  For pure states,
\begin{equation}
 D_{\rm tr}(|\psi\rangle,|\phi\rangle)
 =\sqrt{1-|\braket{\psi}{\phi}|^2}.
 \label{eq:pure-trace}
\end{equation}
Here and below, ket arguments of $D_{\rm tr}$ are shorthand for the
corresponding rank-one projectors.
This follows by diagonalizing the rank-two Hermitian operator
$|\psi\rangle\langle\psi|-|\phi\rangle\langle\phi|$, whose two nonzero
eigenvalues are
$\pm\sqrt{1-|\langle\psi|\phi\rangle|^2}$.

Using
$\langle0|R_Y(a)^\dagger R_Y(b)|0\rangle=\cos((a-b)/2)$ gives the exact
distance
\begin{equation}
 d_N(c,c')
 :=D_{\rm tr}(|\Phi(c)\rangle,|\Phi(c')\rangle)
 =\sqrt{1-\prod_{j=1}^{N}
 \cos^2\!\left(\frac{\alpha_j(c)-\alpha_j(c')}{2}\right)}.
 \label{eq:product-distance}
\end{equation}

For a broadcast encoder $\alpha_j(c)=\alpha(c)$,
\begin{equation}
 d_N(c,c')=\sqrt{1-\cos^{2N}\!\left(
 \frac{\alpha(c)-\alpha(c')}{2}\right)}.
 \label{eq:broadcast-distance}
\end{equation}
At every fixed $N$, inputs with
$|\alpha(c)-\alpha(c')|<\pi$ are nonorthogonal, so
Corollary~\ref{cor:orthogonal} rules out exact conditional universality.

Assume now that $\alpha$ is $L_\alpha$-Lipschitz.  Write
$u=\cos^2(\Delta\alpha/2)$.  The elementary inequalities
$1-u^N\leq N(1-u)$ and $|\sin x|\leq|x|$ imply
\begin{align}
 d_N(c,c')
 &=\sqrt{1-u^N}
 \leq\sqrt{N(1-u)}\\
 &=\sqrt{N}\left|\sin\frac{\Delta\alpha}{2}\right|
 \leq\frac{\sqrt{N}L_\alpha}{2}|c-c'|.
 \label{eq:broadcast-lipschitz}
\end{align}

\begin{corollary}[Sharp-transition counterexample]
\label{cor:sharp-ramp}
Let a continuous target change linearly from $\delta_{0^n}$ to
$\delta_{1^n}$ across a condition interval of width $h$.  Every fixed channel
fed by the $N$-qubit broadcast encoder satisfies
\begin{equation}
 \sup_c\TVD\bigl(q_\theta(\cdot\mid c),p(\cdot\mid c)\bigr)
 \geq\frac12\left(1-\frac{\sqrt{N}L_\alpha h}{2}\right)_+.
 \label{eq:sharp-ramp-bound}
\end{equation}
\end{corollary}

\begin{proof}
The target laws at the two endpoints have total-variation distance one.
Combine Eq.~\eqref{eq:broadcast-lipschitz} with
Theorem~\ref{thm:pairwise}.
\end{proof}

For the reported Gaussian experiment,
$N=19$ and $\alpha(c)=\pi(0.35-0.70c)$, so
$L_\alpha=0.70\pi$.  Taking $h=0.01$ in
Eq.~\eqref{eq:sharp-ramp-bound} gives
\begin{equation}
 \frac12\left(1-\frac{\sqrt{19}(0.70\pi)(0.01)}{2}\right)
 \approx 0.4760.
 \label{eq:paper-numerical-bound}
\end{equation}
This bound holds even if the IQP diagonal is replaced by an arbitrary quantum
circuit with unlimited trainable parameters.

For fixed $c\neq c'$ with
$0<|\Delta\alpha|<\pi$, Eq.~\eqref{eq:broadcast-distance} approaches one as
$N\to\infty$.  Repetition therefore evades a fixed pairwise bound only by
making the product states asymptotically orthogonal.  Simultaneously,
Eq.~\eqref{eq:broadcast-lipschitz} shows that the best common local modulus
scales as $\sqrt{N}$.  This is the label-amplification limit identified in
Theorem~\ref{thm:asymptotic-separation}.

\section{Two independent fixed-architecture obstructions}
\label{sec:fixed}

The contraction theorem applies to arbitrary trainable channels.  The IQP
and finite-parameter structure impose additional restrictions on the fixed
architectures tested in the paper.

\subsection{Finite Fourier bandwidth}

For the broadcast state in Eq.~\eqref{eq:product-encoder}, the one-qubit
density operator is
\begin{equation}
 R_Y(\alpha)|0\rangle\langle0|R_Y(\alpha)^\dagger
 =\frac12\left(I+\sin(\alpha)X+\cos(\alpha)Z\right).
 \label{eq:rotation-bloch}
\end{equation}
Every entry of the $N$-qubit product density operator is therefore a
trigonometric polynomial of degree at most $N$.  For any fixed POVM effect
$M_A$,
\begin{equation}
 q_\theta(A\mid\alpha)=
 \Tr\!\left(M_A|\Phi(\alpha)\rangle\langle\Phi(\alpha)|\right)
 \label{eq:event-linear-state}
\end{equation}
is therefore a real trigonometric polynomial of degree at most $N$.  This is
the standard Fourier-spectrum restriction induced by repeated angle
encoding~\citep{schuldEffectDataEncoding2021,casasMultidimensionalFourier2023}.

\begin{proposition}[Explicit Fourier counterexample]
\label{prop:fourier-counterexample}
Fix $N$ broadcast input rotations and let $0<\eta\leq1$.  On a full
$2\pi$ period, define a binary-supported target by
\begin{equation}
 p(0^n\mid\alpha)=\frac{1+\eta\cos((N+1)\alpha)}{2},
 \qquad
 p(1^n\mid\alpha)=1-p(0^n\mid\alpha).
 \label{eq:high-frequency-target}
\end{equation}
Then every fixed downstream quantum channel obeys
\begin{equation}
 \sup_\alpha\TVD\bigl(q_\theta(\cdot\mid\alpha),
 p(\cdot\mid\alpha)\bigr)\geq\frac{\eta}{4}.
 \label{eq:fourier-bound}
\end{equation}
\end{proposition}

\begin{proof}
Let $g(\alpha)=q_\theta(0^n\mid\alpha)$.  It has zero Fourier coefficient at
frequency $N+1$, whereas the coefficient of
$f(\alpha)=p(0^n\mid\alpha)$ at that frequency is $\eta/4$.  If
$\lVert g-f\rVert_\infty<\eta/4$, then the magnitude of every Fourier
coefficient of $g-f$ is smaller than $\eta/4$, contradicting its coefficient
$-\eta/4$ at frequency $N+1$.  Hence
$\lVert g-f\rVert_\infty\geq\eta/4$.  Total variation dominates the
probability difference of the event $\{0^n\}$, which proves
Eq.~\eqref{eq:fourier-bound}.
\end{proof}

The full-period assumption is needed only for the explicit constant
$\eta/4$.  On any closed interval $I$ of nonzero length, the restrictions of
degree-$N$ trigonometric polynomials form a finite-dimensional, hence closed,
proper subspace $\mathcal{T}_N$ of $C(I)$.  The function in
Eq.~\eqref{eq:high-frequency-target} is not in that subspace: equality on an
interval would, by analyticity, imply equality on the full real line and
contradict its nonzero $(N+1)$st Fourier coefficient.  Consequently,
\begin{equation}
 \delta_{N,I}:=\inf_{g\in\mathcal{T}_N}
 \left\|g-\frac{1+\eta\cos((N+1)\,\cdot)}{2}\right\|_{\infty,I}>0.
 \label{eq:interval-fourier-gap}
\end{equation}
The same binary-supported target therefore has approximation error at least
$\delta_{N,I}$ on the paper's nondegenerate affine angle interval.  Thus
fixed bandwidth is not dense there either, although
Eq.~\eqref{eq:fourier-bound} gives the simpler numerical bound only on a full
period.

\subsection{Parameter dimension across finitely many conditions}

Fix $K$ encoded inputs and an $n$-bit visible output.  Remove one redundant
probability from each simplex, so the space of arbitrary conditional tables
has dimension
\begin{equation}
 d=K(2^n-1).
 \label{eq:table-dimension}
\end{equation}
An IQP ansatz with $M$ periodic phase parameters defines a real-analytic map
\begin{equation}
 F:\mathbb{T}^M\longrightarrow\R^d.
 \label{eq:analytic-map}
\end{equation}

\begin{proposition}[Dimension obstruction]
\label{prop:dimension}
If $M<d$, the image of $F$ has Lebesgue measure zero in $\R^d$ and is not
dense in the product of $K$ probability simplices.
\end{proposition}

\begin{proof}
The derivative of $F$ has rank at most $M<d$ everywhere, so every point of
$\mathbb{T}^M$ is critical.  Sard's theorem implies that $F(\mathbb{T}^M)$
has $d$-dimensional Lebesgue measure zero.  Because $\mathbb{T}^M$ is compact
and $F$ is continuous, the image is compact and therefore closed.  A closed
measure-zero subset cannot be dense in a simplex product with nonempty
relative interior.
\end{proof}

For a complete two-local diagonal on $Q$ qubits with one one-body and one
two-body phase for each support,
\begin{equation}
 M=Q+\binom{Q}{2}.
 \label{eq:two-local-parameters}
\end{equation}
The reported settings give the following necessary-dimension failures.  The
Gaussian row uses its five-condition training grid; the Curie--Weiss row uses
its six training temperatures.

\begin{center}
\begin{tabular}{lrrrr}
\toprule
Experiment & visible $n$ & conditions $K$ & phases $M$ & $K(2^n-1)$\\
\midrule
Gaussian & 19 & 5 & 190 & 2,621,435\\
Curie--Weiss & 8 & 6 & 91 & 1,530\\
\bottomrule
\end{tabular}
\end{center}

This count concerns exact coverage or density of arbitrary tables at the
listed fixed inputs.  It does not claim that the experimental targets require
all table dimensions, nor does it replace the pairwise contraction bound.

\section{Classical tractability is orthogonal to universality}
\label{sec:ct}

The paper's classical training method~\citep{recioArmengolTrainClassical2025}
requires input states for which one can
sample computational-basis strings and evaluate the amplitude ratios entering
Pauli estimators.  Product rotation states satisfy these requirements.  So do
constant product states, which give the immediate counterexample after
Corollary~\ref{cor:orthogonal}, and orthogonal computational-basis labels,
which enable routing.  Tractability therefore includes examples at both
extremes and does not control conditional expressivity.

More formally, tractability is an algorithmic property of a state
description.  The inequalities in Sections~\ref{sec:contraction} and
\ref{sec:resource} depend only on the density operators themselves.  No
variance reduction, amplitude-ratio identity, or classical sampling oracle
can increase $D_{\rm tr}(\rho_c,\rho_{c'})$, because the physical output is
still obtained by the channel in Eq.~\eqref{eq:povm-channel}.

Conversely, orthogonal selector states may be trivially classically
tractable, while a generic condition-addressed diagonal can require an
exponential number of phases to describe.  Thus three resource questions must
be kept separate:
\begin{enumerate}
  \item whether the input admits the paper's classical estimator;
  \item whether condition states remain overlapping or become routing labels;
  \item whether the diagonal has a polynomial-size local description.
\end{enumerate}

Recent unconditional IQP universality results concern the ability of a
sufficiently general diagonal to reproduce one target law, sometimes with
hidden qubits~\citep{kurkinUniversalityKernelAdaptive2025}.  Compiling many
such laws behind a selector is a useful construction, but
Proposition~\ref{prop:selector-mixture} shows why it does not establish the
shared conditional universality studied here.

\section{Admissible positive statements and open boundary}
\label{sec:boundary}

The no-go result suggests the strongest possible target class for a fixed
encoder.  If a model approximates $p$ with uniform error $\varepsilon$, then
Theorem~\ref{thm:pairwise} necessarily gives, for all $c,c'$,
\begin{equation}
 \TVD\bigl(p(\cdot\mid c),p(\cdot\mid c')\bigr)
 \leq D_{\rm tr}(\rho_c,\rho_{c'})+2\varepsilon.
 \label{eq:compatible-kernel}
\end{equation}
Therefore, a defensible positive theorem must restrict targets to kernels
whose total-variation geometry is dominated by the trace-distance geometry of
the chosen encoding.  Even Eq.~\eqref{eq:compatible-kernel} is only necessary,
not sufficient: locality, Fourier bandwidth, hidden width, and parameter
sharing impose additional constraints.

For the complete two-local IQP family, it remains open whether a useful
nontrivial subclass of encoding-compatible kernels admits a uniform
approximation theorem with polynomial resources.  Such a theorem would be a
genuinely shared conditional result.  It would have to exploit relations
between nearby conditions rather than assign them to orthogonal blocks.

\section{Conclusion}
\label{sec:conclusion}

The requested broad, non-multiplexed universality statement is false.  The
proof is the channel-contraction inequality
\begin{equation*}
 \TVD(q_\theta(\cdot\mid c),q_\theta(\cdot\mid c'))
 \leq D_{\rm tr}(\rho_c,\rho_{c'}),
\end{equation*}
which yields an explicit approximation lower bound for every pair of
conditions.  Exact universality requires pairwise orthogonal condition
states.  Uniform universality for all continuous conditional laws requires
asymptotic pairwise orthogonalization and the loss of every resource-uniform
continuity modulus.  These states are operational condition labels, so this
route to broad universality implements the selector or routing behavior
excluded in the question.

Discarded selectors do not provide a coherent loophole: an IQP diagonal turns
them into the classical mixture in Eq.~\eqref{eq:selector-mixture-law}.  The
paper's fixed broadcast product encoder additionally admits the quantitative
counterexample in Eq.~\eqref{eq:sharp-ramp-bound}, while finite Fourier support
and insufficient parameter dimension independently rule out universality of
the tested fixed-width complete two-local circuits.  The appropriate claim is
therefore a no-go theorem, together with the encoding-compatibility condition
in Eq.~\eqref{eq:compatible-kernel}, rather than a selector-based universality
construction.

\bibliographystyle{unsrtnat}
\bibliography{paper/supplement_theory}

\end{document}
